\scanpage{13}

par composition \inserted{de foncteurs} entre catégories satisfaisant \unsure{un|aux} ens. également donné de conditions d'exactitude. Pour ceci, il est \deleted{\unsure{bien|bon}} naturel d'exiger que

\diagram{\begin{tikzcd}[ampersand replacement=\&,row sep=1.4em,column sep=1.1em]
\&C\arrow[r,dash,"F"]\&C'\arrow[r,dash,"G"]\&C''\\
\&\Theta\arrow[u,"u"']\&\bar J'\arrow[l,"\eta'"']\&{}\arrow[d,phantom,"{\bigg\}\ \in\lambda=(\lambda',\lambda'')}" description]\\
\bar I'\arrow[ur,"\xi'"]\&\bar I''\arrow[u,"\xi''"']\&\bar J''\arrow[ul,"\eta''"]\&{}\\
{}\arrow[r,phantom,"{\underbrace{\hspace{3.3em}}_{\in\delta=(\delta',\delta'')}}" description]\&{}\&\&
\end{tikzcd}}

\deleted{\unsure{les foncteurs sont les|les foncteurs}} $\lambda'\supset\delta',\lambda''\supset\delta''$. Si on veut que les foncteurs identiques satisfassent ces conditions, il est naturel d'exiger qu'on se borne aux $C$ qui satisfont la condition d'exactitude que tt $\Theta\to C$ $\delta$ exact est $\lambda$ exact. Mais alors la condition \deleted{sur} $F$ devient que $\forall\,\Theta\overset{u}{\unsure{-|\to}} C$ $\lambda$-exact, $F\circ u$ est $\lambda$-exact.

\deleted{On a ainsi donné des moyens de définir}

\deleted{$1^\circ$) \unsure{De} \inserted{un ens $\Sigma_0$} de parties \inserted{$\Delta_0$} de $\mathrm{Ob}(\mathrm{Cat})$, chacune stable par équivalence de catégories (si $C\in\Delta_0$, et $C'\in\mathrm{Ob}\,\mathrm{Cat}$, $C'\cong C$, alors $C'\in\Delta_0$) $\Sigma$ est stable par intersection.}

\deleted{$2^\circ$) Un ens $\Sigma_1$ de parties $\Delta_1$ de $\mathrm{Fl}(\mathrm{Cat})$, chaque $\Delta_1$ stable par isom (si $u,v:C\rightrightarrows C'$, et $u\simeq v$, alors $u\in\Delta_1\Rightarrow v\in\Delta_1$) \uline{contenant les} \unsure{identi|identité} équivalences et par composition}

\textcircled{1} Donnée d'un paquet \inserted{$\rho$} de catégories $\Theta$, avec $\lambda=(\lambda',\lambda'')$. Ici $\rho$ désigne \deleted{\unsure{les}} l'ens des couples $(\Theta,\lambda)$

\deleted{$\Delta_\rho\subset(\mathrm{Cat})$ : les catégories qui}

$\Sigma_\rho\subset\mathrm{Fl}(\mathrm{Cat})$ : un foncteur $F:C-C'$ est $\rho$-exact ssi pour tt \inserted{$(\Theta,\lambda)$ de $\rho$ et tt} foncteur $u:\Theta\to C$ qui est $\lambda$-exact, $F\circ u$ est $\lambda$-exact.

Alors \quad a) $\Sigma_\rho$ est stable par isom (i.e. si $F,G:C-C'$, \deleted{\unsure{et}} $F\in\Sigma_\rho\Rightarrow G\in\unsure{\Sigma|\xi}_\rho$)

b) $\Sigma_\rho$ contient les équivalences de catégories

c) $\Sigma_\rho$ stable par composition

d) Une intersection de parties $\Sigma_\rho$ de $\mathrm{Fl}\,\mathrm{Cat}$ est de $=$ type.

\textcircled{2} Donnée d'un \deleted{\unsure{d'un paquet|de foncteurs}}\inserted{un couple $\sigma\ \deleted{\unsure{\text{de}}}\ (\sigma_1,\sigma_2)$} de paquets de foncteurs $R\xrightarrow{\varphi}\bar R$, avec \deleted{\unsure{$\varphi$}} $R$ muni de $\delta=(\delta',\delta'')$, $\bar R$ de $\bar\delta=(\bar\delta',\bar\delta'')$, et \deleted{\unsure{tels et que}} l'on définit $\Delta_\sigma\subset\mathrm{Ob}\,\mathrm{Cat}$. par

% archive sheet number: 12
